%!TEX root = ../../main.tex
%% ===============================================================
%% 9.2 표현 효과와 학습기 효과의 분해
%% 파일: chapters/09_results/02_decomposition.tex
%% 상위: chapters/09_results/chapter.tex
%% 정의 라벨: sec:results-decomp
%% 참조 라벨: tab:baseline
%% 포함 객체: -
%% 인용: grinsztajn2022why
%% 상태: 초안 (docs/OUTLINE.md 참고)
%% ===============================================================

\section{표현 효과와 학습기 효과의 분해}
\label{sec:results-decomp}

RQ2와 RQ3는 표 \ref{tab:baseline}의 두 간격으로 답할 수 있다.

\begin{itemize}[leftmargin=1.5em]
  \item \textbf{표현 효과 (RQ2).} 같은 신경망 학습기 계열 안에서 테이블 요약(MLP, $\aucMLP{}$)과 최선의 비테이블 표현(BiGRU·융합, $\aucBiGRU{}$)의 차이는 $\gapRepresentation{}$이다. 30 초 관측 창이 60 초 프레임 하나 또는 둘만 참조하는 상황에서, 시퀀스·그래프 구조는 계단 유지된 거의 동일한 상태를 여섯 번 반복해서 보는 셈이고, 실제로 변하는 것은 희소한 사건 계수뿐이다. 반면 테이블 요약 $\phi_{\mathrm{tab}}$은 그 희소한 변화를 명시적 통계량(차분, 기울기)으로 노출한다.
  \item \textbf{학습기 효과 (RQ3).} 입력을 고정한 LightGBM($\aucLGBM{}$)과 MLP($\aucMLP{}$)의 차이 $\gapLearner{}$는 순수한 학습기 효과이다. 트리 앙상블은 축 정렬 분할과 잎당 최소 표본 제약으로 수천 차원의 이질적 특징에서 상호작용을 안정적으로 찾는 반면, MLP는 같은 표본 수에서 과적합과 최적화의 어려움을 겪는다. 이는 테이블 데이터에서 그래디언트 부스팅이 심층 신경망보다 우세하다는 광범위한 보고 \cite{grinsztajn2022why}와 일치한다.
\end{itemize}

두 효과의 크기가 비슷하다는 점은 중요하다. 신경망 표현이 뒤처진 이유의 절반은 표현 자체가 아니라 학습기 계열에 있으며, 분 단위 텔레메트리에서는 어떤 표현을 쓰든 \emph{드러낼 수 있는 신호의 총량}이 상한을 이룬다는 해석을 지지한다.
